\section{表头类别图例、连续色阶与共享图例}
\subsection{横向类别图例与组合表头}
\key{\gradtablelegend[选项]{名称列表}} 按自然宽度横向排列类别色块与文字；
\key{\gradtablelegendhead[选项]{标题}{名称列表}} 在其上方居中放置普通文本标题，间隔 2pt。
两者均属于核心，\key{libraries=none} 即可使用，无需加载 \key{legend} 库。
它们只在显式调用时显示，不会自动插入其他图形或表格。
\begin{example}[text and listing,sidebyside=false]{服务响应时间：标题与五项状态图例}
\begingroup
\gradbarssetup{stack colors={gray,teal,blue,
  orange,red}}
\begin{center}
\begin{tabular}{lc}
服务 & \gradtablelegendhead{响应状态}{待测,快速,正常,偏慢,超时}\\
\hline
查询 & \gradbar[max=1000,width=42mm,color=teal,precision=0,unit={ ms}]{120}\\
导出 & \gradbar[max=1000,width=42mm,color=blue,precision=0,unit={ ms}]{420}\\
\end{tabular}
\par\medskip
\gradtablelegendhead[table legend=false]{响应状态}{待测,快速,正常,偏慢,超时}
\end{center}
\endgroup
\end{example}
这些是 \key{/gradbars} 图形键，可放在助手选项或 \key{\gradbarssetup} 中，
遵循普通 TeX 分组与从左到右覆盖；不是 \key{gradtable} 的新增表格声明键。
\key{table legend=true|false} 默认 true，裸键等同 true；false 隐藏类别行，
标题助手仍显示标题，且不保留 2pt 图例间距。
\key{legend font} 默认 \key{\sffamily\scriptsize}，仅设置类别文字；
\key{legend swatch size=6pt} 必须为正，
\key{legend text gap=2pt} 与 \key{legend item gap=4pt} 必须非负。
边框默认为 \key{legend border color=black,legend border width=.4pt}，线宽必须非负。

类别复用 \key{\gradbarscategory} 的名称映射，未注册名称按 \key{stack colors} 的顺序循环取色；
打印模式沿用类别纹理。图例不会推断阈值、给观测分类或更改条形颜色，示例因此显式设置条形颜色。
已有竖向 \key{\gradbarslegend} 和堆叠条的 \key{legend=true} 不受 \key{table legend} 控制。
标题助手可直接用于普通表头或 \key{\gradcolumn} 的标题参数；收集数值列使用：
\begin{codeonly}
\gradbarshead{\gradtablelegendhead{响应状态}
  {待测,快速,正常,偏慢,超时}}
\end{codeonly}
可用该助手替换原先手写的标题、minipage 与 TikZ 色块宏；助手本身不插入 \key{\multicolumn}，
跨列仍由外层表格代码决定。完整独立示例见 \key{docs/examples/service-response-legend.tex}。
类别图形沿用 \key{alt}、\key{accessibility} 与 \key{category-legend} 后端规则，
需要原生回退时仍须显式选择；标题是 Figure 之外的普通文本。

\subsection{连续色阶、共享图例与参考标记}
导言区加载 \key{legend} 库；其依赖 Heat。连续图例表达数值尺度，复用 Heat 的插值函数，
不能把类别色块当成连续量程。
\begin{example}[text and listing,sidebyside=false]{两格数值共享冻结的尺度与图例}
\begingroup
\gradscale{legend-score}{min=-10,max=30}
\gradlegenddefine{legend-score}{scale group=legend-score,
  heat palette=purple-teal,
  heat center=0,legend title={Score}}
\gradheat[legend group=legend-score]{-5}
\gradheat[legend group=legend-score]{20}
\par\medskip
\gradlegenduse[width=60mm,height=3mm,legend missing=false]{legend-score}
\endgroup
\end{example}
\key{legend missing=false} 隐藏缺失值色块及其描述，默认值为 \key{true}；两个图例命令均可使用。
\key{\gradheatlegend} 接受普通图形选项，以及 \key{legend title} 和 \key{legend ticks}。\par
方向由 \key{legend orientation} 设置，取 \key{horizontal} 或 \key{vertical}。
width 表示色带长度，height 表示厚度；默认标注两端点和内部发散中心。
自定义刻度必须在量程内严格递增，位置按原值计算，文字按精度与单位格式化。
色带采用 256 段 Heat 插值底色，并叠加装饰性半透明白色分隔线；
水平、纵向和打印模式均保留这些细纹。它们改变外观，不改变数值尺度，也不表示类别或分箱边界。
刻度过密或舍入后相同，应增加精度或调整布局。
\key{\gradlegenddefine} 在当前分组中冻结显式尺度、方向、单位和颜色角色；
独立图形、原生列的 options 和 tabularray 单元格均可引用 \key{legend group=name}。
未知名称、重复定义及尺度／配色／单位冲突会报错。
\key{\gradlegenduse} 仅允许布局、刻度格式、字体与替代文本覆盖，不改变后续图形设置。
\begin{example}[text and listing,sidebyside=false]{参考符号与统计含义}
\gradreferencelegend{target/Declared target,
  reference/Reference value,band/Reference range,
  error/Error distances,box/Quartiles and observed whiskers}
\end{example}
文字必须说明标记的真实含义：误差距离不自动等于置信区间，箱线须也不是误差须。
共享类别或参考图例可用 \key{legend kind=categorical|reference} 与 \key{legend names}。
重复装饰图例可设 \key{accessibility=artifact}，完整图例的自动描述保留尺度、中心、方向与单位。


\subsection{内置 Heat 配色}
用一行 \key{heat palette=名称} 同时设置两端颜色、白色中心和插值模式。
顺序配色为 \key{blue}、\key{teal}、\key{purple}、\key{orange}，从白色渐变到对应颜色；
发散配色为 \key{purple-teal}、\key{blue-red}、\key{orange-blue}、\key{brown-teal}，
颜色顺序对应低值、白色中心、高值。发散模式仍由 \key{heat center} 指定中心数值。
这些名称是配色快捷方式，不等于可访问性认证；原有 accessible 配色继续保留。
预设之后仍可用 \key{heat low color}、\key{heat center color}、\key{heat high color} 单独覆盖。
\begin{example}[text and listing,sidebyside=false]{内置连续色阶配色}
\gradheatlegend[width=45mm,height=2mm,min=-10,max=30,heat center=0,
  heat palette=blue,legend title={blue},legend missing=false]\quad
\gradheatlegend[width=45mm,height=2mm,min=-10,max=30,heat center=0,
  heat palette=teal,legend title={teal},legend missing=false]\par\medskip
\gradheatlegend[width=45mm,height=2mm,min=-10,max=30,heat center=0,
  heat palette=purple,legend title={purple},legend missing=false]\quad
\gradheatlegend[width=45mm,height=2mm,min=-10,max=30,heat center=0,
  heat palette=orange,legend title={orange},legend missing=false]\par\medskip
\gradheatlegend[width=45mm,height=2mm,min=-10,max=30,heat center=0,
  heat palette=purple-teal,legend title={purple-teal},legend missing=false]\quad
\gradheatlegend[width=45mm,height=2mm,min=-10,max=30,heat center=0,
  heat palette=blue-red,legend title={blue-red},legend missing=false]\par\medskip
\gradheatlegend[width=45mm,height=2mm,min=-10,max=30,heat center=0,
  heat palette=orange-blue,legend title={orange-blue},legend missing=false]\quad
\gradheatlegend[width=45mm,height=2mm,min=-10,max=30,heat center=0,
  heat palette=brown-teal,legend title={brown-teal},legend missing=false]
\end{example}
